\chapter{工作总结与展望}
本章对全文工作进行概括，并说明后续可以继续推进的方向。工作总结部分围绕语义建模、系统设计与验证、原型实现与实验评估三个方面展开；未来展望部分则从形式化验证、工程实现、事务协调、实验评测和自动配置等角度说明当前工作的延伸空间。

\section{工作总结}

本文围绕“分布式数据库如何在保证事务语义正确的前提下支持混合隔离级别，并提升整体执行效率”这一问题展开研究。现有工作往往分别停留在语义分析、形式化验证或系统实现中的某一环节，三者之间缺少稳定衔接。针对这一不足，本文尝试将语义定义、操作语义设计、形式化验证、原型实现和实验评估串联起来，形成一条从理论约束到系统验证的研究链路。

在语义建模方面，本文基于抽象执行模型，对事务历史、可见性关系 $\text{VIS}$、仲裁关系 $\text{AR}$ 以及一致性公理进行了统一描述，并给出了 RA、CC、PC、PSI、SI 和 SER 等隔离级别的表达方式。对于混合隔离级别场景，本文强调系统正确性不能被理解为各事务局部正确性的简单叠加，而需要在统一抽象执行下进行全局判定。

在系统设计与验证方面，本文面向分布式键值数据库给出了一套与抽象执行语义相对应的操作语义。该设计以客户端、路由节点、配置节点和分片节点组成的分布式架构为基础，引入 HLC、MVCC、依赖追踪、等待机制、检查集维护、冲突检测和两阶段提交等机制，将不同隔离级别的可见性与顺序约束落实到可执行的状态转移规则中。考虑到 \TLAPlus 适合描述并穷举检查状态转移系统，本文进一步借助 TLC 对关键安全性质进行了模型检查，在给定有限模型范围内确认系统状态演化不会产生违反目标一致性约束的执行。

在工程实现与实验评估方面，本文设计并实现了一个支持 PC、PSI、SI 和 SER 混合运行的分布式键值数据库原型系统。该原型采用 Router--Shard 架构，集成了分片管理、路由转发、HLC 时间戳维护、MVCC 存储、多分片事务协调和单分片快速提交等核心机制。基于 Courseware、SmallBank 和 TPCC 三组负载，本文在单可用区集中部署和跨可用区分散部署两种拓扑下进行了对比测试。实验结果表明，相较于统一采用强隔离级别的基线配置，混合隔离级别执行机制在本文测试条件下能够提高系统吞吐量；扩展性对照实验也显示，更细粒度的事务分配策略仍有进一步改善空间。

综上，本文的主要工作包括三个方面：建立混合隔离级别场景下的一致性语义与抽象执行判定框架；设计并验证面向分布式键值数据库的混合隔离级别操作语义；通过原型系统和实验评测说明该机制在给定实验条件下具备工程可实现性和一定性能收益。

\section{未来展望}

本文工作仍带有明显的研究型原型特征，后续可以从以下几个方向继续推进。

首先，可以扩展形式化建模与验证范围。本文的 \TLAPlus 模型主要关注事务可见性、提交顺序和隔离级别约束，对故障恢复、副本复制、网络分区和节点动态变化等因素进行了抽象。后续可在保持模型可分析性的前提下，引入故障场景、恢复过程与副本一致性约束，进一步判断混合隔离级别机制在复杂分布式环境中的正确性边界。

其次，可以完善原型系统的工程能力。当前系统仍以研究验证为目标，对持久化、日志恢复、复制容错、在线重配置和弹性伸缩等机制支持有限。后续可优先补齐存储引擎、日志子系统和复制协议，并评估这些底层机制对混合隔离级别执行路径的影响。

再次，可以优化跨分片事务协调机制。本文主要依赖两阶段提交协议保证跨分片原子性，该方案实现清晰、便于分析，但在高并发、广域网络或参与分片较多的场景下会放大协调开销。后续可研究更高效的原子提交机制，或探索与复制协议、时间戳定序机制结合更紧密的事务协调方法。

此外，可以扩大实验评测范围。后续可引入更标准化和更大规模的基准测试，进一步分析不同读写比例、事务长度、热点分布和跨分片比例下的系统表现，并补充长尾延迟、资源利用率、事务中止率等指标，以更全面地评估混合隔离级别机制的正确性与性能权衡。

最后，可以进一步研究隔离级别推荐与自动配置问题。本文主要关注固定隔离级别集合下的统一语义与系统支持，而真实业务中的隔离需求往往与操作模式、冲突特征和业务约束共同相关。若能将自动配置策略与本文的统一语义框架结合，系统将不仅支持多种隔离级别共存，还能更有针对性地为不同事务选择执行策略。
